\documentclass[11pt]{letter}
\usepackage[utf8]{inputenc}
\usepackage[T1]{fontenc}
\usepackage{geometry}
\geometry{margin=1in}
\usepackage{hyperref}
\hypersetup{colorlinks=true,linkcolor=blue,urlcolor=blue,citecolor=blue}
\usepackage{enumitem}

\pagestyle{empty}

\begin{document}

\begin{flushright}
March 12, 2026
\end{flushright}

\bigskip
\noindent
Editor-in-Chief\\
\textit{npj Wireless Technology}

\bigskip
\noindent
Dear Prof.\ Alouini,

\medskip
\noindent
We are pleased to submit our manuscript entitled: \textbf{``Truly Mobile Sub-Terahertz Communications Beyond 6G via Just-in-Time IMU-Assisted Beam Management''} for consideration in \textit{npj Wireless Technology}, Collection: \textit{Extreme Band Communications: Innovations in THz and Optical Wireless Technologies} (\url{https://www.nature.com/collections/hbfcgcbggg})

\medskip
\noindent
Sub-THz frequencies are a candidate enabling technology for 6G-and-beyond wireless networks. While stationary sub-THz links have demonstrated impressive data rates, maintaining narrow ``pencil beam'' alignment under realistic user mobility (combining translation and unpredictable rotation) remains an open challenge. Existing reactive beam management strategies, even when augmented with inertial measurement unit (IMU) data, are insufficient to prevent frequent outages and are thus not suitable for truly mobile links. In addition, the lack of end-to-end sub-THz evaluation hardware platforms that jointly accounts for high-mobility dynamics, real-time sensor data, and real-life characteristics of sub-THz hardware limits progress in this important area.

\medskip
\noindent
To address this, we present: (i)~\emph{Just-in-Time (JIT)}, a hybrid proactive plus reactive intelligent beam management algorithm that fuses RF and IMU data to preemptively realign the link \emph{before} failure; and (ii)~\emph{MobileTHz}, an integrated simulation and Hardware-in-the-Loop (HIL) platform for mobile sub-THz systems. Our numerical study concludes that \emph{JIT delivers \underline{up to 1--2} \underline{orders-of-magnitude improvements in link availability} over state-of-the-art approaches while incurring \underline{less than 0.4\% signaling overhead}}. Simulation results are also cross-validated through HIL measurements with sub-THz RF hardware operating in the 110--170\,GHz band.

\medskip
\noindent
\textbf{Why \textit{npj Wireless Technology} and the ``Extreme Band Communications: Innovations in THz and Optical Wireless Technologies'' collection?}\\
Our work addresses a fundamental bottleneck in mobile (sub-)THz communications, efficient beam management under nodes mobility. Solving this task in an essential building block on the way from research and prototyping extreme-band communication systems toward their future practical deployments. To our knowledge, this is one of the first comprehensive studies of a hybrid proactive+reactive beam realignment technique for mobile sub-THz validated through HIL measurements. The open-source MobileTHz platform further provides a reusable tool for the community.\\We believe this fits well within the collection's scope on THz innovations for 6G and beyond.

\medskip
\noindent
This manuscript has not been previously published and is not under consideration elsewhere.\\All authors have approved the submission. We declare no competing interests.

\medskip
\noindent
\textbf{Recommended referees}
\begin{enumerate}[leftmargin=1.5em, itemsep=2pt]
    \item Prof.\ Marco Giordani (University of Padova, Italy) \href{mailto:giordani@dei.unipd.it}{giordani@dei.unipd.it}
    \item Prof.\ Yasaman Ghasempour (Princeton University, USA) \href{mailto:ghasempour@princeton.edu}{ghasempour@princeton.edu}
    \item Prof.\ Zhi Sun (Tsinghua University, China) \href{mailto:zhisun@tsinghua.edu.cn}{zhisun@tsinghua.edu.cn}
    \item Prof.\ Robert W.\ Heath Jr.\ (UC San Diego, USA) \href{mailto:rwheathjr@ucsd.edu}{rwheathjr@ucsd.edu}
\end{enumerate}

\medskip
\noindent
Thank you very much for considering our manuscript.

\bigskip
\noindent
Sincerely,

\noindent
Sergey Petrushkevich\\
Ph.D.\ Candidate\\
Department of Electrical and Computer Engineering\\
\& Institute for the Wireless Internet of Things\\
Northeastern University\\
Boston, MA 02115, USA\\
Email: \href{mailto:petrushkevich.s@northeastern.edu}{petrushkevich.s@northeastern.edu}

\end{document}